\section{Introduction}
A finite memory in an experiment must preserve the prediction directions
that the past has actually acquired and the future can still distinguish.
It must also permit every subsequent update without recovering a
discarded command. These requirements are not determined by the
ambient dimension of a posterior family. When an experiment degenerates,
different prediction directions can disappear at different rates, and
their attainable dimension can be smaller than the future test space.

The principal result of this paper identifies this geometry uniformly
through arbitrary additive exponent collisions in positive monomial
experiments. A family of exterior volumes gives the optimal checkpoint
error for every integer memory budget. The same family, maximized over
checkpoints, gives the error of one causal filter. The constants remain
uniform at the collision itself. The proof connects three distinct
assertions: unconditional mass acquired by positive histories, an
anisotropic cover of the entire reachable image, and stable updates of
reachable representatives in coordinates with no reciprocal collision
gaps.

We place this result within a structural theory for command-polynomial
experiments on arbitrary compact latent spaces. That theory separates
rank from resolution, and the dependence on a specified prior from
attainment under every full-support prior. Positive affine acquisition
provides an exact covariance model for the resolution geometry. A
quotient of the product space distinguishes prescribed annihilation
from exact kernel realization. The circular detector, the inverse of
the checkpoint budget curve, and prior ambiguity then describe other
aspects of the same acquired--observed distinction. Each conclusion
below is stated with its own uniformity and resource assumptions.

\subsection{The collision-uniform memory law}

For the monomial experiment write
\[
 A=\{0=a_0<a_1<\cdots<a_{r-1}=D\},\qquad
 W_A=\operatorname{span}\{t^a:a\in A\},\qquad h_A(m)=|mA|,
\]
where $mA$ consists of sums of exactly $m$ members of $A$, including
zero entries. A positive finite detector spans $W_A$. A command in
$[\eta,1-\eta]^J$, with $0<\eta<1/2$, retains or rejects each cell;
rejection is a report, and every trial is counted. All report
likelihoods are at least $\kappa>0$. The exact information statements
hold on $I=[l,u]$, $0\le l<u<\infty$, with a full-support prior.
For the uniform collision statements take $I=[0,1]$, fix $N\ge2$,
and let the one-step exponents vary over a compact subset $K$ of the
strictly ordered chamber. One fixed full-column-rank coefficient matrix
specifies the positive detector on $K$, and $\mu$ is fixed with full
support; it need not have a density.

Choose $H\ge\max\{1,N\max_K a_{r-1}\}$. Keep every nonconstant
formal future label $\alpha\in\N_0^r$, $|\alpha|=m$, even when
exponents coincide. Put
\[
 q_m=\binom{m+r-1}{r-1}-1,\qquad x_\alpha=H^{-1}\alpha\cdot a,
 \qquad
 \mathcal V_{m,\ell}(a)=\max_{|F|=\ell}
       \prod_{\{\alpha,\beta\}\subset F}|x_\alpha-x_\beta|.
\]
Thus $\mathcal V_{m,1}=1$; a zero product remains zero. Set
\[
 \begin{split}
 p_{n,m}&=\min\{n(r-1),q_m\},\\
 \Psi_{n,m}(M,a)&=\max_{1\le\ell\le p_{n,m}}
       (\mathcal V_{m,\ell}(a)/M)^{2/\ell},\\
 \Xi_N(M,a)&=\max_{1\le n<N}\Psi_{n,N-n}(M,a).
 \end{split}
\]
Endpoint profiles are zero. The checkpoint query is a uniformly selected
member of a fixed finite spanning menu of physical $m$-trial products,
selected independently after the memory index is formed. Only the
selected query is executed. Commands during acquisition are independent
uniform commands for the unconditional criterion. A checkpoint code may read the complete acquired prefix before retaining
one of $M$ labels. A causal code must update that label after each report
without reading discarded commands. Regret is the excess squared loss
for the selected query; its unconditional and worst-history versions
are compared below. Persistent labels, the read-only clock, and the fixed
program are specified in Definition~\ref{def:finite-state}.

\begin{theorem}[Intrinsic attainable resolution and causal memory]\label{thm:resolution-main}
Under these hypotheses there are constants $0<c\le C<\infty$,
depending only on the fixed experiment, $K$, $N$ and $\mu$, with the
following properties for every known calibration $a\in K$ and every
integer $M\ge1$.

At each checkpoint $n+m=N$, the optimal $M$-message squared-prediction
regret is between $c\Psi_{n,m}(M,a)$ and $C\Psi_{n,m}(M,a)$,
both unconditionally under independent uniform acquisition commands and
in the minimax sense over admitted histories. The query is sampled
independently after the index is formed, from the ordered products of
one fixed attainable failure-factor basis; all its $m$ trials are counted.

There exists one deterministic stage-compatible transducer, allowed to
depend on $a$ and $M$, storing at most $M$ persistent labels after
every report, whose regret at every checkpoint and on every history is
at most $C\Xi_N(M,a)$. Conversely every such transducer has maximum
unconditional checkpoint regret at least $c\Xi_N(M,a)$ under the
same common exploration law. Independent coding randomization does not
remove the lower bound. The clock and fixed real-valued program are
read-only as in Definition~\ref{def:finite-state}.

The constants are uniform across all additive collision strata of $K$.
No calibration-blind codebook or uniformity over all full-support priors
is part of the assertion.
\end{theorem}


The exact continuous-state dimension is separately
$d_A(n,m)=\min\{n(r-1),|mA|-1\}$, as stated in
Theorem~\ref{thm:main}. The finite-resolution formula specifies how
these directions disappear. Its proof joins three assertions not
contained in an ambient spectral estimate. Actual interior failure
factors have a separated binomial product tangent. Pairing it with
complete Newton--Hermite prefixes loses exactly the evidence direction
under normalization and supplies a positive-history minorization.
The entire posterior image has bounded semialgebraic format, so its
thin-rectangle covering bound truncates at the same attainable
dimension. Finally, raw-moment updates have a Lipschitz recurrence
without inverse collision gaps, and quantization of reachable
representatives realizes one causal filter. These are established in
Lemmas~\ref{lem:newton-attainment} and \ref{lem:tame-rectangle} and
Theorems~\ref{thm:intrinsic-checkpoint} and
\ref{thm:intrinsic-streaming}.


\subsection{From attained directions to one causal filter}
The scales in the preceding theorem belong to the interaction of past
and future. A future Vandermonde determinant alone supplies neither
history mass nor a streaming algorithm. The proof first pairs the
actual binomial product tangent with complete Newton--Hermite prefixes.
After removal of evidence, every attainable prefix has uniformly
positive unconditional mass. The past truncation $n(r-1)$ therefore
enters the lower bound before quantization.

For the upper bound, every report-word posterior image is a rational
image of its command cube. Bounded semialgebraic format and its
attainable dimension control a cover of the whole image in a thin
rectangle. This step does not require a single global regular chart.
For the causal part, the representative is a reachable raw-moment
vector. Multiplication by the next likelihood closes the remaining test
space, and its positive evidence denominator bounds the update on
posterior-mixture segments. The finite-horizon error recurrence includes
every preceding quantization error. Thus the theorem joins local
acquisition, global geometry, and temporal compatibility, rather than
using any one of them as a substitute for the others.

The structural and affine results below explain which parts of this
argument are general and which depend on monomial acquisition. In
particular, a dimension law has no collision-uniform constants unless
additional geometry is proved, while an affine covariance law does not
by itself supply the multi-step Newton flags.

\subsection{Structural rank and prior strata}

Let $X$ be a compact metric space, let $\mu$ be a probability on $X$,
and fix a finite horizon $N$. At time $i$ the command belongs to a
finite-dimensional cube, the report alphabet is finite, and
\[
 L_{i,y}(u,x)=\sum_{\alpha\in F_{i,y}}u^\alpha c_{i,y,\alpha}(x),
 \qquad c_{i,y,\alpha}\in C(X),\qquad L_{i,y}\ge\kappa>0,
 \qquad\sum_yL_{i,y}=1.
\]
The polynomial condition concerns $u$, not $x$. In particular, finite
detectors with affine retention gates belong to this class for arbitrary
continuous detector cells. Future word likelihoods span a
finite-dimensional test space. Choose a finite spanning list of actual
queries, each with a fixed positive selection probability, and write
$p_n(h)$ for the weighted vector of their conditional success
probabilities after a history of length $n$.

For each past report word $a$, let $P_a$ be its likelihood product as a
function of the commands, set $Z_a=\mu P_a$, and put
\[
 Y_a=(Z_a,Z_ap_n)^t,\qquad
 d_n=\max_a\rank_{\R(\boldsymbol u)}[Y_a,DY_a]-1.
\]
All entries are command polynomials with finitely many actual prior
moments as coefficients. The evidence column is adjoined before
normalization; it is not counted as an acquired prediction direction.
Theorem~\ref{thm:v18-rank-classification} proves, whenever $d_n>0$,
\[
 R^{\rm av}_{M,n}\asymp R^{\rm max}_{M,n}\asymp M^{-2/d_n},
 \qquad
 \mathfrak R^{\rm av}_{M,N}\asymp
 \mathfrak R^{\rm max}_{M,N}\asymp M^{-2/d_*},
 \quad d_* =\max_{1\le n<N}d_n>0.
\]
Here $R$ is checkpoint regret and $\mathfrak R$ is the largest
checkpoint regret of one causal filter. The average criterion uses
independent uniform commands and includes report evidence; the maximum
criterion ranges over all admitted histories. A zero rank gives a
finite reachable set, not necessarily a singleton. If all checkpoint
ranks vanish, a finite exact causal filter exists.

This formula classifies each specified experiment without assuming an
accessible flag or an ordered-scale representation of its risk.
Theorem~\ref{thm:v18-moment-strata} further identifies every locus
$d_n\le k$ by the vanishing of finitely many polynomials in the prior
moments: they are the command coefficients of the $(k+2)$-minors of the
augmented matrices. Full-support priors fill the relative interior of the finite
moment body. Simultaneous maximal ranks hold off a proper algebraic
set there, and can be reached by arbitrarily small bounded-density
perturbations of any full-support prior. Exceptional priors are not
excluded from the classification; their actual ranks enter the formula.
Uniform constants hold on compact constant-rank moment sets. A rank
formula by itself does not assert uniformity across changing ranks.

\subsection{Affine acquisition and realizable degeneration}

The first full-class resolution theorem applies to positive affine
acquisition. Let $f\in C(X)^b$ with
$\sup_x\sum_i|f_i(x)|\le a<1$, choose a uniform command
$u\in[-1,1]^b$, and observe a report $y\in\{-1,1\}$ with likelihood
\[
 L_y(u,x)=\tfrac12(1+y u^tf(x)).
\]
For any finite menu of future query probabilities, let $h$ be its vector
weighted by the square roots of the query probabilities and define
\[
 C_\mu=\Cov_\mu(h,f).
\]
The physical command cube fixes the input units. No singular whitening
is performed. If $s_1\ge\cdots\ge s_d\ge0$ are the singular values
of $C_\mu$, Theorem~\ref{thm:v18-affine-classification} proves
\[
 R_M^{\rm av}\asymp R_M^{\rm max}\asymp
 \max_{1\le\ell\le d}
      \left(\frac{s_1\cdots s_\ell}{M}\right)^{2/\ell}.
\]
The constants depend only on $a,b$ and the number of queries. They are
uniform over priors, query weights, multiplicities, and all rank losses.
The assertion includes the one-label exact case $C_\mu=0$.

The proof derives this geometry before addressing quantization. With
$v=yu$ and $m=\mu f$, the exact prediction vector is
\[
 p=\mu h+C_\mu\theta,\qquad
 \theta=\frac{v}{1+m^tv}.
\]
The law of $v$, including both reports and their evidence, has density
$(1+m^tv)/2^b$. Its projective image contains a fixed ball, is contained
in another fixed ball, and has a uniformly positive density on the
inner ball. Thus the same covariance image controls both global
approximation and unconditional acquisition mass. This is not an
inference from an ambient observation matrix. As a concrete interior
rank change, acquisition through $x$ and prediction through $x^2$ under
$d\mu_t=(1+tx)\,dx/2$ on $[-1,1]$ have covariance proportional to
$t$. Proposition~\ref{prop:v18-interior-example} obtains the full
uniform $t^2M^{-2}$ law, including zero at $t=0$, while every prior
still has full support.

A second conclusion concerns rank degeneration inside a dominated
prior family. For each pair $b,q$, one fixed positive acquisition and
one fixed physical query menu admit priors $\mu_A$ with
\[
 \tfrac12\mu_0\le\mu_A\le\tfrac32\mu_0,\qquad
 C_{\mu_A}=\frac{A}{8b\sqrt q},\qquad
 \sum_{j,i}|A_{ji}|\le\tfrac12.
\]
Thus arbitrary small matrix germs, and not only a scalar cancellation,
are realized by genuine priors in one experiment
(Theorem~\ref{thm:v19-covariance-realization}). Along an analytic
arc let $A_\ell$ be the least vanishing order among its $\ell$-minors.
Theorem~\ref{thm:v19-analytic-memory} gives, uniformly in small
$t>0$ and all integer $M\ge1$,
\[
 R_{M,1}(t)\asymp
     \max_{1\le\ell\le r}t^{2A_\ell/\ell}M^{-2/\ell}.
\]
The differences $\alpha_\ell=A_\ell-A_{\ell-1}$ are ordered
nonnegative integers. Every such finite list occurs under the same
physical bounds, and the full checkpoint budget curve determines the
$A_\ell$ operationally. In particular the phase exponent along
$M=\lceil t^{-\gamma}\rceil$ is
$\min_\ell 2(A_\ell+\gamma)/\ell$. This is a complete arcwise
classification for affine acquisition, not a simultaneous normal form
for all polynomial history maps in several parameters.

\subsection{Universal attainment and exact kernel constraints}
Theorem~\ref{thm:v18-universal-pairing} addresses a different
quantifier. At an interior positive history let
$\widetilde T=\operatorname{span}\{P\}+\operatorname{im}DP$ and
$\Phi=\operatorname{span}\{1,\phi_1,\ldots,\phi_p\}$. When both
spaces have dimension $p+1$, normalized acquisition has rank $p$ for
\emph{every} full-support prior if and only if
\[
 \det[v_i(x_k)]_{i,k=0}^p\det[\phi_j(x_k)]_{j,k=0}^p
\]
has one weak sign on $X^{p+1}$ and is nonzero somewhere. Here $v$ is
any basis of $\widetilde T$. Necessity follows by approximating atomic
probabilities by full-support probabilities; sufficiency follows from
determinant integration. Separate Chebyshev properties are sufficient
instances, not necessary hypotheses. In particular,
$\widetilde T=P\Phi$ gives prior-uniform full rank on arbitrary compact
latent spaces. The square hypothesis is explicit; no reduction of every
rectangular pairing to this criterion is asserted.


For finite-dimensional $E,F\subset C(X;\R)$ with $\dim F\le\dim E$,
Theorem~\ref{thm:v19-rectangular} characterizes full pairing rank
under every full-support prior by
\[
 0\ne f\in F\quad\Longrightarrow\quad
 \text{there exists }g\in E\text{ with }fg\ge0,\quad fg\not\equiv0.
\]
Theorem~\ref{thm:v19-rank-alternative} characterizes subspaces that
can be contained in a pairing kernel, and constructs full-support priors
attaining the minimum rank. Exact realization of a specified subspace
requires a further condition.
Adjoining the positive history product before normalization converts
pairing rank into actual prediction rank. This yields the unconditional
mass conclusion without a square reduction. The five-point experiment
in Proposition~\ref{prop:v19-pentagon} satisfies the rectangular
criterion but has no fixed universally nonsingular two-dimensional
tangent subspace containing evidence.


The exact-kernel question admits a finite multiplication criterion.
For a specified $U\subset F$ set
\[
 W_U=\spanop(EU),\qquad V=\spanop(\{1\}\cup EF).
\]
Annihilation is feasible under full support precisely when
$W_U\cap C(X)_+=\{0\}$. Choose a basis of $V/W_U$ beginning with
$[1]$, and express the pairing on $E\times(F/U)$ in this basis. This
gives an affine matrix pencil $\mathscr H_U(z)$ depending only on
pointwise multiplication modulo $W_U$. Theorem~\ref{thm:exact-kernel}
proves
\[
 \exists\mu\text{ of full support with }\ker H_\mu=U
 \quad\Longleftrightarrow\quad
 W_U\cap C(X)_+=\{0\},\quad
 \rank_{\R(z)}\mathscr H_U=\dim(F/U).
\]
The determinantal condition is a polynomial identity test on the full
affine quotient, not an existential restatement on an unknown prior
slice. Positive density perturbations give local coordinates on that
slice and attain its maximal rank arbitrarily close to every feasible
prior. The common forced kernel is
$\{f\in F:Ef\subset W_U\}$, which need not equal $U$.
Proposition~\ref{prop:block-kernels} computes both the forced kernel
and the exact criterion in a family of positive block experiments by
the common zeros of a polynomial subspace. This includes examples where
annihilation is feasible but exact realization fails, and explicit
full-support witnesses for the subspaces that are exactly realizable.
Corollary~\ref{cor:kernel-history-rank} returns this distinction to the
actual normalized prediction rank and its unconditional history mass.

\subsection{Contrast, inversion, and unresolved prior information}

A circular experiment with Haar prior has the positive cells
$\tfrac14(1\pm\tau\cos\theta)$ and
$\tfrac14(1\pm\tau\sin\theta)$. Its four lookup probabilities vary
continuously in $[\eta,1-\eta]^4$. For every fixed horizon and all
sufficiently small known $\tau\ge0$, the checkpoint law is
\[
 H_{n,m}(M,\tau)=
   \max_{1\le j\le\min(n,m)}\tau^{2(j+1)}M^{-1/j}.
\]
The $j$th harmonic is attenuated by $\tau^j$ when acquired and again
by $\tau^j$ when queried. The real physical scales are therefore
$\tau^2,\tau^2,\tau^4,\tau^4,\ldots$, not just a count of Fourier
coefficients. Theorem~\ref{thm:circle-resolution} proves a uniform
causal realization through weighted Fourier updates without inverse
contrast. The contrast interval and constants belong to the fixed
horizon, not to an unbounded-horizon assertion.

Within these established profile families the complete checkpoint curve
has an inverse. For $R(M)$ define
\[
 \mathcal I_\ell(R)=\inf_{M\in\N,\,M\ge1}M R(M)^{\ell/2}.
\]
Theorem~\ref{thm:operational-reconstruction} and
Corollary~\ref{cor:v18-affine-invariant} identify this transform with
the initial scale products up to uniform constants. Integer budgets
require a factor-two sandwich; the underlying real-budget identity is
exact. Repeated scales and zero products are included. This inverse
uses the entire unbounded budget curve, not finitely observed errors,
and does not recover a latent model or exact optimal constants. Nor is
it applied to a maximum over checkpoints or to a curve with a positive
uncertainty floor.

Prior uncertainty has its own geometry. For monomials, the local
ambiguity body after a positive history has successive width orders
$\varepsilon d_1,\ldots,\varepsilon d_{q_m}$, where $d_i$ are the
future Leja scales. There is no past-dimension truncation of this
ambiguity body. After an exact prefix of $k$ raw future moments,
the remaining squared radius is $\asymp\varepsilon^2d_{k+1}^2$.
The normalization of prior tilts is explicit in
Proposition~\ref{prop:ambiguity-ellipsoid}. For an imperfect common
moment name of error $\delta$, the separate dominated-family and
interior-slack hypotheses of Theorem~\ref{thm:sharp-common-moments}
give the sharp joint law $\Xi_N(M,a)+\delta^2$. The lower bound
uses fixed genuine prior alternatives and overlapping history laws,
not alternatives chosen after seeing a history.


\subsection{Classical inputs and the scope of the conclusions}
The analytic inputs include semialgebraic dimension and metric entropy,
determinant integration, total positivity, and divided differences
\cite{BochnakCosteRoy,ForresterAndreief,Pinkus,deBoor,
YomdinComte,ZhangKileel,ComteHalupczok}. The finite square-pairing
criterion is the strong-compatibility result of Banaji--Pantea
\cite{BanajiPantea2016}; the finite rectangular sign-vector alternative
is present in M\"uller et al.\ \cite{MullerEtAl2016}. The exact
specializations are recorded in Section~\ref{sec:v19-comparison}.
The product-character expansion used to perturb the covariance is
classical Fourier analysis on the finite cube
\cite[Section 2.2]{deWolf2008}. Analytic invariant orders and their
singular-value interpretation are classical as well
\cite{KavehMakhnatch}. We do not claim any of these inputs as an
independent advance in their respective theories.

What must additionally be proved in the experiment is the passage from
these inputs to attainable prediction error with the stated quantifiers.
For monomials this is the collision-uniform conjunction of positive
history mass, the dimension-truncated whole-image cover, and one
stage-compatible causal filter. For affine acquisition, Bayes' formula
gives the global covariance image and its actual probability law before
quantization. The analytic-arc law is a consequence of that physical
classification and classical matrix orders, not a separate normal-form
theory. The kernel theorem supplies the missing exact-realization
condition together with positive prior witnesses; the underlying
separation and determinantal arguments remain elementary.

\begin{table}[tbp]
\centering\footnotesize
\setlength{\tabcolsep}{4pt}
\begin{tabular}{@{}p{.16\textwidth}p{.23\textwidth}p{.34\textwidth}p{.21\textwidth}@{}}
\toprule
Result & Acquisition and latent space & Uniformity & Memory conclusion\\
\midrule
Collision law & Monomial, multiple steps; interval & Fixed full-support
prior; fixed horizon and compact ordered chamber; all additive collisions
& All integer budgets; checkpoint and one causal filter\\[3pt]
Structural rank & Command-polynomial; arbitrary compact space & Every
specified prior; uniform on compact constant-rank moment sets & Rank
exponent; checkpoint and one causal filter\\[3pt]
Affine profile & One affine step; arbitrary compact space & Across all
covariance rank losses, for fixed physical bounds and dimensions & Full
checkpoint profile; not a general multi-step profile\\[3pt]
Exact kernels & Continuous finite product spaces & Full-support
annihilation slice; arbitrarily small bounded positive tilts & Exact local
rank and acquired mass; global and causal claims require more\\[3pt]
Circular law & Multiple steps; circle and Haar prior & Fixed horizon and
known sufficiently small contrast, including zero & Paired scales;
checkpoint and one causal filter\\[3pt]
Prior ambiguity & Monomial moment uncertainty & Stated domination and
interior-slack hypotheses & Widths; joint memory and common-moment error\\
\bottomrule
\end{tabular}
\caption{Hypotheses and conclusions. In every memory statement the query
is chosen after the label is formed. Persistent labels are charged;
known calibration, clock, and the fixed program follow the stated
read-only convention. No discarded command tape is free.}
\label{tab:scope}
\end{table}

The manuscript first develops the monomial experiment, its attainable
transversality, and the collision-uniform law. It then gives the general
rank and affine classifications, the square, rectangular, and exact
kernel criteria, and the associated positive-history and analytic
consequences. The circular classification, operational inverse, and
prior-uncertainty geometry follow. The observation algebra, confluent
specializations, sequential decision results, and numerical constructions
appear with their proofs in the appendices.

The separate finite implementation theorem,
Theorem~\ref{thm:effective-main}, supplies integer transition tables and
dyadic outputs under its explicit numerical interface. Advice, program
length, transient workspace, and persistent labels are distinct resources.
Sharpness is asserted for persistent memory; sufficient bounds are
supplied for the other resources. The intrinsic results do not presume
computability of unspecified real prior moments. Detailed comparisons
with finite control approximations and information-state methods are in
Appendix~\ref{sec:literature-comparison}.
